Attention matrices are often said to collapse in rank during training, and spectral statistics are attractive as cheap guides for pruning. Most evidence, however, comes from finished models. We test five specific claims: heads split into high- and low-rank populations, entropy follows layer-specific shapes, heads below $0.3\log n$ entropy can be pruned safely, $I(A;X)\le H(Y|X)+O(1/\sqrt n)$, and the last layer reaches its rank minimum three to five times faster than the first. We train 68 transformers across 13 configurations on text, vision, and pretrained \bert, tracking every head at 41 checkpoints, and run 120 further models with controlled $H(Y|X)$. No run meets the bifurcation criterion under any of five rank definitions. Under default initialization attention starts nearly uniform and rank rises. With label signal, rank falls substantially only when attention starts peaked, and then the last layer is 1.8 times faster than the first (95\% CI 1.6--2.0), not three to five times. The entropy threshold removes 63--96\% of heads on \agnews and \cifar at a cost of 19--45 percentage points. We give a counterexample to the information bound and prove a valid replacement based on attention variance. Rank collapse is therefore a property of the initial state, not a general law of training. Entropy helps rank heads for pruning on vision tasks, but a fixed threshold is unsafe.
